\documentclass[11pt,a4paper]{article}
\usepackage{amsmath,amssymb}
\begin{document}
\title{How to Recognize Classical Chaos\\and What Quantum Spectra Can (and Cannot) Tell Us}
\author{}
\date{}
\maketitle

\section*{What is being characterized?}
There is no single test that captures every use of the word ``chaos.'' An orbit may be sensitive to perturbations without a chosen invariant measure having positive entropy; a system can mix without being uniformly hyperbolic; and a spectral plot is not itself a classical trajectory test. It helps to separate three questions: how nearby states move, how much information a typical orbit produces, and which features a particular experiment can resolve. We first consider a differentiable map $T:M\to M$ on a compact phase space with an invariant probability measure $\mu$. For a flow $\Phi^t$, the analogous statements use time in place of iteration, or a specified return map. The measure, energy surface, and time unit must be stated when quoting an entropy or exponent.

\section*{Geometric and trajectory-based viewpoints}
\paragraph{Lyapunov exponents.} A perturbation $v$ at $x$ evolves to $DT_x^n v$. When its long-time rate exists, its Lyapunov exponent is
\[\lambda(x,v)=\lim_{n\to\infty}\frac{1}{n}\log\frac{\|DT_x^n v\|}{\|v\|}.\]
Positive values express exponential linearized separation in some directions, not exponential separation of two bounded trajectories forever. The linear approximation fails once their separation is macroscopic. Under the integrability assumptions of the multiplicative ergodic theorem, tangent space splits into subspaces with well-defined exponents for $\mu$-almost every point; for an ergodic measure their values are almost everywhere constant. A finite-time estimate depends on the orbit, direction, norm, and observation window. A chaotic sea can coexist with regular islands, so one should not replace a measure-average by one conspicuously unstable trajectory. For a volume-preserving invertible map, the sum of the exponents, counted with multiplicity, is zero. An autonomous Hamiltonian flow normally has neutral directions associated with the flow and conserved quantities; its exponents should be computed on the appropriate energy shell or transverse section.

\paragraph{Hyperbolicity, stretching, and folding.} Uniform hyperbolicity is a strong geometric condition: tangent directions split continuously into uniformly contracting and expanding bundles (with a neutral flow direction for a flow). It supplies a robust mathematical mechanism for sensitive dependence, but many systems called chaotic are not uniformly hyperbolic. Stretching by itself is not enough to characterize recurrent bounded dynamics: folding, global returns, invariant sets, and coexistence with stable regions matter. A Poincar\'e section can expose these structures for a flow. Intersections of stable and unstable manifolds, horseshoes, and transverse homoclinic orbits provide more specific geometric evidence; a visually complicated phase portrait alone does not establish any of them.

\paragraph{Integrability and regular motion.} A Hamiltonian with sufficiently many independent commuting integrals can, under the regularity and compactness hypotheses of the action--angle construction, move quasiperiodically on invariant tori. Typical torus orbits have zero Lyapunov exponents. The hypotheses and the word ``typical'' matter: singular separatrices or unstable special orbits can occur even in otherwise integrable systems. Conversely, failing to find a conserved quantity is not a proof of chaos. In two degrees of freedom, sections displaying invariant curves alongside scattered-looking regions suggest mixed dynamics, but rigorous conclusions require more than the image.

\section*{Information and statistical viewpoints}
\paragraph{Kolmogorov--Sinai entropy.} Let $\mathcal P$ be a finite measurable partition and let $\mathcal P^{(n)}=\bigvee_{k=0}^{n-1}T^{-k}\mathcal P$ distinguish length-$n$ symbolic histories. With Shannon entropy $H_\mu(\mathcal P^{(n)})=-\sum_{A\in\mathcal P^{(n)}}\mu(A)\log\mu(A)$, define
\[h_\mu(T,\mathcal P)=\lim_{n\to\infty}\frac{H_\mu(\mathcal P^{(n)})}{n},\qquad h_\mu(T)=\sup_{\mathcal P}h_\mu(T,\mathcal P).\]
This is an asymptotic information-production rate for an invariant measure, not the entropy of an instantaneous probability density. Positive $h_\mu$ means that some finite-resolution observation produces new information per iteration; it need not say that every observable does. A generating partition, when available under the relevant hypotheses, permits the supremum to be evaluated from one partition. Changing the invariant measure can change the KS entropy of the same map.

\paragraph{When do entropy and exponents agree?} Pesin's entropy formula, in settings where its differentiability and invariant-measure hypotheses hold (notably a $C^{1+\alpha}$ diffeomorphism of a compact smooth manifold with an absolutely continuous invariant probability measure), states
\[h_\mu(T)=\int_M\sum_{\lambda_i(x)>0}m_i(x)\lambda_i(x)\,d\mu(x),\]
where $m_i$ is the multiplicity. The integral, rather than an arbitrary single-orbit sum, is essential for a nonergodic measure. General relationships between metric entropy and positive exponents require care with smoothness and measure assumptions; the displayed equality is not a blanket identity for every singular invariant measure. Its left side concerns the entire dynamical system and an asymptotic supremum over partitions. A finite-time slope of a coarse-grained physical entropy, including observational entropy, is not automatically that left side.

\paragraph{Ergodicity, mixing, and correlations.} Ergodicity says that invariant measurable sets have measure zero or one. Mixing is stronger: for measurable $A,B$, one has $\mu(A\cap T^{-n}B)\to\mu(A)\mu(B)$. For observables, this appears as relaxation of correlations $\int f\,(g\circ T^n)\,d\mu-\int f\,d\mu\int g\,d\mu$. Exponential, algebraic, or absent correlation decay distinguishes different dynamical regimes, but depends on the chosen class of observables. Ergodicity and mixing are not interchangeable with a positive numerical Lyapunov exponent or a positive KS entropy. Stronger properties, including the measure-theoretic K-property and Bernoulli behavior, describe increasingly restrictive forms of statistical unpredictability. Statistical tests over a short finite record cannot by themselves establish those asymptotic properties.

\paragraph{Topological entropy.} In contrast to $h_\mu$, topological entropy $h_{\mathrm{top}}(T)$ counts the exponential growth of distinguishable orbit segments at arbitrarily fine resolution without first choosing an invariant probability measure. For a continuous map of a compact metric space, the variational principle reads
\[h_{\mathrm{top}}(T)=\sup_{\mu\,\mathrm{invariant}}h_\mu(T).\]
Thus a map can have rich topological orbit complexity while a measure supported on a periodic orbit has zero metric entropy. Topological entropy, KS entropy, and a measured entropy slope answer different questions.

\paragraph{Recurrences and transport.} Poincar\'e recurrence on a finite-measure invariant region is generic and hence not, by itself, a chaos test. Return-time statistics, escape rates in an open system, diffusion of a transported coordinate, and decay of survival probabilities can be informative, but their form also depends on the region, observable, and invariant structures. Apparent diffusion or broad-band power spectra can be mimicked by noise or by unresolved multiscale regular motion. These are useful diagnostics to combine with tangent-space and phase-space evidence, not universal definitions.

\section*{A simple calibration example}
For the area-preserving torus map $T(x)=Ax\pmod{1}$ with $A=\begin{pmatrix}2&1\\1&1\end{pmatrix}$, the expanding eigenvalue is $(3+\sqrt5)/2$. Its positive Lyapunov exponent with respect to invariant area is $\log((3+\sqrt5)/2)$ per iteration; in this smooth hyperbolic example the KS entropy of invariant area has the same value. The agreement has an identifiable mechanism and measure. By contrast, an invariant measure concentrated on the fixed point still has zero KS entropy even though the ambient map is expanding and contracting nearby. This illustrates why ``the exponent'' and ``the entropy'' should always be accompanied by an orbit ensemble or invariant measure.

\section*{Bethe ansatz and energy statistics belong to the quantum side}
\paragraph{Bethe ansatz.} The Bethe ansatz is a technique for constructing eigenstates and energies of particular quantum many-body Hamiltonians. When applicable and complete, it provides evidence for a special kind of quantum integrability. It is not a characterization of classical chaos. Some integrable models are not solved by this method, and failure to find a Bethe-ansatz solution is not evidence that a model is chaotic. Even identifying quantum integrability by conserved charges requires checking independence, locality or other intended restrictions, and the relevant sector; simply naming a solvable limit does not classify its classical dynamics.

\paragraph{Eigenenergies versus level correlations.} A raw eigenenergy histogram mainly describes the density of states and can conceal both symmetries and scale changes. Quantum spectral diagnostics instead resolve symmetry sectors, unfold local levels where appropriate, and examine gaps $\delta_n=E_{n+1}-E_n$. For unfolded mean spacing one, $P(s)=e^{-s}$ is the Poisson benchmark often associated with integrable spectra; the familiar $P(s)=(\pi/2)s\exp(-\pi s^2/4)$ is a \emph{Wigner surmise}, not the exact spacing law of every time-reversal-invariant chaotic system. Symmetry class, degeneracies, integrability conditions, and spectral window all matter. The adjacent-gap ratio $r_n=\min(\delta_n,\delta_{n-1})/\max(\delta_n,\delta_{n-1})$ is another useful statistic that does not require unfolding. Neither a single histogram nor level repulsion proves classical chaos: the correspondence between classical dynamics and quantum spectral statistics is semiclassical and conditional, not a definition. For a periodically driven system one examines Floquet eigen\emph{phases} (quasienergies) modulo $2\pi$, rather than treating them as an ordinary unbounded energy spectrum.

\paragraph{Operator growth and observational entropy.} Out-of-time-ordered commutators, operator spreading, and observational entropy provide additional quantum diagnostics. A semiclassical commutator can approximate a Poisson bracket over a limited time window, but its finite-time growth is not automatically a classical Lyapunov exponent. Observational entropy depends on the chosen POVM and, in a finite-dimensional system, ultimately has finite capacity. A fitted rise of that entropy need not equal the sum of positive classical exponents: one must specify the semiclassical limit, the measurement resolution, the state ensemble, and the order of time and system-size limits. Likewise, a full measurement-record entropy and a single-time observational entropy must not be conflated with KS entropy.

\section*{Practical reading of a chaos claim}
A careful comparison should report the phase space and invariant measure; whether an exponent is local, finite-time, or asymptotic; the definition of entropy and its partition or POVM; the symmetry sector and unfolding used in spectral plots; and any finite-size or semiclassical extrapolation. The strongest conclusion supported by an isolated calculation is usually that one diagnostic is \emph{consistent with} a proposed chaotic regime. Agreement between a KS entropy and a positive-exponent sum needs a theorem's hypotheses; agreement of a quantum entropy slope with either of them additionally needs a separate, model- and measurement-dependent argument.

\end{document}
